% !TEX root =  master.tex
\chapter{Einleitung}

\section{Problemstellung \& Motivation} 
Mit der fortschreitenden Entwicklung künstlicher Intelligenz (KI) stehen Unternehmen vor der Herausforderung, ihre Geschäftsprozesse kontinuierlich zu optimieren und durch den gezielten Einsatz neuer Technologien effizienter zu gestalten. Insbesondere im Audit-Umfeld bietet der Einsatz von KI ein erhebliches Potenzial zur Steigerung von Effizienz, Qualität und Transparenz. Die Prüfung von Dokumenten und Prozessen muss dabei zeiteffizient, nachvollziehbar und möglichst vollständig durchgeführt werden. KI kann interne Auditoren unterstützen, indem sie den manuellen Analyseaufwand reduziert und gleichzeitig umfassendere sowie schnellere Prüfungen ermöglicht (vgl. \cite[S. 11--12]{wassie2024artificial}).

Ein besonders vielversprechendes Anwendungsfeld stellt die Analyse von Audit-Stichproben dar. Im Rahmen von Prüfungen werden häufig Dokumente wie Rechnungen, User-\allowbreak Provisioning-Vorgänge oder Verträge untersucht, um Auffälligkeiten, Regelverstöße oder potenzielle Risiken zu identifizieren. Die manuelle Analyse dieser Unterlagen ist jedoch
zeitaufwendig und bindet erhebliche personelle Ressourcen. Gleichzeitig steigt aufgrund wachsender Datenmengen und zunehmend komplexer Geschäftsprozesse die Schwierigkeit,
relevante Auffälligkeiten zuverlässig zu erkennen.

Vor diesem Hintergrund gewinnt die Unterstützung von Stichprobenanalysen durch KI-basierte Verfahren zunehmend an Bedeutung. Für die Freudenberg-Gruppe als global agierendes Unternehmen stellt sich dabei die Frage, wie KI-Technologien gewinnbringend in den Internal-Audit-Prozess integriert werden können. Neben der Verbesserung der Effizienz und Prüfungsqualität spielen insbesondere die Einhaltung strenger Datenschutz- und Compliance-Anforderungen sowie die Auswahl eines geeigneten KI-Modells für die Analyse von Audit-Dokumenten eine zentrale Rolle. Daraus ergibt sich die Notwendigkeit, die Eignung verschiedener KI-Lösungen für die Unterstützung von Stichprobenanalysen im Internal Audit systematisch zu untersuchen.

\section{Zielsetzung der Arbeit}
Ziel dieser Projektarbeit ist die systematische Untersuchung der für die Stichprobenanalyse 
\section{Forschungsfrage}
Kann auch in ein Überkapitel - gibt es Unterfragen?
\section{Aufbau der Arbeit}